\subsection{What does showing a full-list partner guarantee?}\label{sec:two-theory}

The full-list partner in this section is an assignment in the same model as the restricted replay. The Chilean historical offer provides a separate empirical target.

\paragraph{The best prediction can send two people to one seat}
Three applicants share three one-seat programs $x,y,z$. The vector of first choices, in applicant order, takes the values $(x,y,z)$, $(x,z,y)$, $(y,x,z)$ and $(y,z,x)$ with probabilities $2/6,2/6,1/6,1/6$. Each state has distinct first choices, so all three receive those choices under full-list DA whatever the program priorities. The platform knows the probabilities but cannot see the state. With one slot per applicant, the unique choices maximizing expected partner coverage are $(x,z,z)$. They cover $5/3$ partners in expectation and place two people. The lists $(x,y,z)$ cover $3/2$ partners in expectation and place all three. Thus filling the third seat can cost as little as $1/6$ of an expected correct recommendation in this example.

\paragraph{What an omitted partner can explain}
Let $H$ be the applicants whose full-list assignment is a real program missing from $D$. Let $K_j$ and $T_j$ count members of $H$ assigned to $j$ under $M$ and $N_D$, respectively.
\begin{theorem}\label{thm:two-omission}
Relative to the full market, gross losses and gross gains in placement on $D$ are each at most $|H|-\sum_j\min(K_j,T_j)$. Furthermore, every unmatched applicant whose full-list partner is displayed can be paired with a distinct member of $H$ who is placed on $D$ and strictly prefers their full-list partner. Both conclusions hold for arbitrary stable selections on the two graphs.
\end{theorem}
\begin{proof}
Delete $H$ together with the seats its members occupy in each assignment. Any remaining blocking pair would have blocked the original assignment. Restrict the residual full-market allocation to the residual display: every surviving contract is visible there, and removing other edges cannot create a block. The two residual allocations are thus stable on the same pairs, using capacities $c-K$ and $c-T$, where $c$ is the original capacity vector. For fixed pairs, adding one seat retains every matched applicant and adds at most one: regard seats as unit copies with common order within each program, run seat-proposing DA, and admit the new proposer last. An applicant who holds an offer can only trade up. Matched-set invariance transfers this comparison to any stable selection \citep{galesoto,roth1986,crawford1991}. Comparing the two capacity vectors through their coordinatewise minimum bounds losses among the remaining applicants by $\sum_j(T_j-K_j)_+$ and gains by $\sum_j(K_j-T_j)_+$. Exactly $|H|-\sum_jT_j$ of $H$ lose a place, and none gain. Adding gives the bounds.

For the individual claim, at each program cancel the holders common to both assignments. Pair departing holders with incoming holders one-to-one until one side is exhausted. Start with an applicant who is unmatched on $D$ despite seeing their full-list partner $j$. Program $j$ must be full in $N_D$, and every new holder there outranks this applicant. Follow the pairing at $j$ to one of its new holders $h$. Full-list stability implies that $h$ had an assignment strictly preferred to $j$, since $j$ previously held a lower-priority applicant. If $h$ still sees that former assignment, repeat. Each step starts with someone who strictly prefers their former assignment and each pairing is injective. On a continued step the program is full after restriction. Its intake of new holders is therefore at least its number of departing holders, so the departing applicant always has a successor. The starting applicant has no incoming link, so the paths cannot cycle or meet. Each path ends at a placed applicant whose former assignment is hidden.
\end{proof}

For two policies $A,B$, one entry in the empirical comparison counts applicants with $C_i(B)-C_i(A)=1$ and $P_i(B)-P_i(A)=-1$. The theorem gives each such applicant a distinct other person on the final display who is assigned a worse program than their full-list partner, which is hidden under $B$. That other applicant can have the same list under $A$ and $B$. It can happen with only two applicants and one slot each. Let $i$ prefer $x$ to $y$, let $h$ prefer $y$ to $x$, and let both programs give $h$ higher priority. Full-list partners are $x$ for $i$ and $y$ for $h$. Display $A$ offers $i$ only $y$ and $h$ only $x$. Both place. Change only $i$'s list to $x$. Now $i$ sees their full-list partner and loses the competition for it. Under one-slot lists every witness is this direct: a new holder at the focal program could not also see their own strictly preferred full-list partner. The traced links in a larger market describe the two assignments and need not follow the chronology of DA proposals.

Let $p_{ij}(X)$ be the probability that $j$ is the full-list partner of $i$, given the information $X$ used to choose the lists. If a policy also uses randomness independent of the full allocation conditional on $X$, condition on its seed as well. By linearity, the conditional expected number covered is $\sum_i\sum_{j\in D(i)}p_{ij}(X)$. Selecting the most probable partners separately for each person therefore minimizes expected omissions, the coarse upper bound $|H|$ above. A smaller upper bound can coexist with fewer actual placements. For comparison, \citet[Theorem~4.7]{pruned} identify which hospitals can become unmatched under one-sided prefix pruning. Here omissions can occur anywhere in an incomplete list, and a displaced person can still see their own full-list partner.
